\documentclass[11pt]{article}
\usepackage{amsmath,amssymb,amsthm,mathtools,mathrsfs}
\usepackage[margin=1in]{geometry}
\usepackage{hyperref}
\newtheorem{theorem}{Theorem}
\newtheorem{proposition}[theorem]{Proposition}
\newtheorem{lemma}[theorem]{Lemma}
\newtheorem{definition}[theorem]{Definition}
\DeclareMathOperator{\Tr}{Tr}
\DeclareMathOperator{\coker}{coker}
\DeclareMathOperator{\Spec}{Spec}
\DeclareMathOperator{\rem}{rem}
\DeclareMathOperator{\diag}{diag}
\newcommand{\C}{\mathbb C}
\newcommand{\F}{\mathbb F}
\newcommand{\A}{\mathbb A}
\newcommand{\Gm}{{\mathbb G_m}}
\newcommand{\HH}{\mathcal H}
\newcommand{\BB}{\mathscr B}
\title{The original marked product at its boundary:\\
the full parameter connection, period determinant,\\
and finite-field mixed extension}
\author{Split-Zero research continuation}
\date{13 September 2026}
\begin{document}
\maketitle

\begin{abstract}
For the original monic packet polynomial, retaining all selected zero
multiplicities, we calculate the missing $u$-connection of its polynomial
exponential complex in the original monomial lattice.  The connection has
an explicit second-order coefficient, a first-order coefficient of trace
$d/2$, and a degree-correct chain realization.  We calculate its full
tensor version and the exact period determinant, including its gamma
factors and ordered contour phase.  We also apply Deligne's lissity and
local-monodromy theorems to the actual finite-field family on $u\ne0$,
obtaining a genuine mixed boundary extension there.  The original theta
source, its Taylor unit, and its metric enter through explicit maps.
Neither a Stokes matrix nor an arithmetic uniform four-volume estimate is
asserted.  The calculations identify and begin the intervening boundary
work rather than declaring the marked programme closed.
\end{abstract}

\section{Fixed objects and original source maps}

Let $Z$ be the original nonempty finite set of selected zeros of
$g=2\xi$, with the complete selected orders $m_\rho$, and put
\begin{equation}
 h(s)=\prod_{\rho\in Z}(s-\rho)^{m_\rho}
     =s^d+\sum_{b=0}^{d-1}h_bs^b,
 \qquad d=\sum_{\rho\in Z}m_\rho\ge1.
 \label{eq:h}
\end{equation}
The primitive and its constant are fixed as
\begin{equation}
 \Phi_h(s)=\frac{s^{d+1}}{d+1}
       +\sum_{b=0}^{d-1}\frac{h_bs^{b+1}}{b+1},
 \qquad \Phi_h(0)=0,
 \qquad f_t(s)=\Phi_h(s)-ts.
 \label{eq:phase}
\end{equation}
Retain the original source spaces, Fourier convention, and operators
\begin{align}
 V&=\{\phi\in\mathcal S(\mathbb R)_{\rm ev}:
          \phi(0)=0,\ \int_{\mathbb R}\phi=0\},\nonumber\\
 \BB&=\{F\in C^\infty(\mathbb R_{>0}):
      \sup_{x>0}x^b|D^mF(x)|<\infty
              \text{ for all }b\in\mathbb Z,m\ge0\},\nonumber\\
 D&=-x\partial_x,\qquad
 \Theta\phi(x)=2\sum_{n\ge1}\phi(nx),\qquad
 \mathcal MF(s)=\int_0^\infty F(x)x^s\frac{dx}{x},
 \label{eq:source}\\
 \widehat\phi(\xi)&=\int_{\mathbb R}\phi(x)e^{-2\pi ix\xi}\,dx.
 \nonumber
\end{align}
The seed is
\begin{equation}
 \phi_*=(4\pi^2x^4-6\pi x^2)e^{-\pi x^2},
 \qquad \mathcal M\Theta\phi_*=g.
\end{equation}
The inherited Euler inverses supply the original $F_h\in\BB$ with
$\mathcal MF_h=v_h=g/h$.  Define
\begin{equation}
 E_h=\C[s]/h,\quad A_h=M_s,\quad
 \upsilon_h=j_hv_h\in E_h^\times,\quad
 \alpha_h(P)=P(D)\phi_*,\quad \mathcal T_h(P)=P(D)F_h.
 \label{eq:maps}
\end{equation}
The original identities are
\begin{equation}
 \Theta\alpha_h(P)=\mathcal T_h(hP),\qquad
 J_h\mathcal T_h(P)=\upsilon_h[P]_h.
 \label{eq:source-square}
\end{equation}
Thus $[\C[s]\xrightarrow{h}\C[s]ds]$ maps to
$[V\xrightarrow{\Theta}\BB]$ by $(\alpha_h,\mathcal T_h)$.
Its degree-one map is exactly $\sigma_hM_{\upsilon_h}$ into
$Q=\BB/\Theta V$.  The linear source section is not asserted to be a
ring homomorphism.  Tensoring these chain maps preserves the full factor
$\upsilon_h^{\otimes k}$ and the original ordered exterior signs.

\section{The full lattice and its marked fibre}

For $R_k=\C[u,t_1,\ldots,t_k]$ set
\begin{equation}
 \mathcal C_k^j=R_k[s_1,\ldots,s_k]\otimes
                  \bigwedge^j(ds_1,\ldots,ds_k),\qquad
 \mathcal D=\sum_{i=1}^k L_i\,ds_i\wedge(-),\qquad
 L_i=u\partial_{s_i}+h(s_i)-t_i.
 \label{eq:complex}
\end{equation}
The $L_i$ commute, so the exterior signs give $\mathcal D^2=0$.

\begin{proposition}[The original complete lattice]
The complex \eqref{eq:complex} has cohomology only in degree $k$, with
\begin{equation}
 \HH_k:=H^k(\mathcal C_k)
   =\bigoplus_{0\le a_i<d}R_k
       [s_1^{a_1}\cdots s_k^{a_k}ds_1\wedge\cdots\wedge ds_k].
 \label{eq:lattice}
\end{equation}
Specialization at $u=t_1=\cdots=t_k=0$ identifies this displayed
module with $E_h^{\otimes k}$, retaining every primary jet.
\end{proposition}
\begin{proof}
In one factor $L=u\partial_s+h-t$ increases the degree of a nonzero
polynomial by exactly $d$ and retains its leading coefficient.
Consequently $L$ is injective.  Subtracting $L(as^j)$ at each highest
degree gives a remainder of degree below $d$.  A polynomial of degree
below $d$ in the image of $L$ must be zero by the same degree identity.
This proves the direct module decomposition
\[
 R[s]=L(R[s])\oplus R[s]_{<d},\qquad R=\C[u,t].
\]
The inverse of $L$ on its image is $R$-linear and contracts the image
two-term complex.  Tensor the direct decompositions in the specified
factor order.  Every tensor summand with a contractible factor contracts
by that factor's homotopy with the sign contributed by preceding degrees;
the cross terms cancel by the tensor differential rule.  Only the top
remainder summand remains.  The splitting is over $R_k$, so it remains a
splitting after the stated coefficient specialization.
\end{proof}

The family is a literal deformation of the original finite polynomial
source images: with $V_* =\alpha_h(\C[s])$ and
$\BB_h=\mathcal T_h(\C[s])$, its one-factor differential is
\begin{equation}
 d^\theta_{u,t}=\Theta|_{V_*}
       +u\mathcal T_h\partial_s\alpha_h^{-1}
       -t\mathcal T_h\alpha_h^{-1}.
\end{equation}
The inverse maps here are defined on the displayed injective images.
Equation \eqref{eq:source-square} gives
$d^\theta_{u,t}\alpha_h=\mathcal T_hL$.
The derivative term is also the original function-space commutator
\[
 \mathcal T_h(P')=(\log x)P(D)F_h-P(D)((\log x)F_h),
\]
because $[\log x,D]=1$.  In Mellin coordinates this subtraction
retains and cancels the two equal $v_h'P$ terms, leaving $v_hP'$.

If $\chi$ is the minimal polynomial of $S=\sum_i s_i$ on
$E_h^{\otimes k}$, the weighted cyclic injection is
\begin{equation}
 \eta:\C[S]/\chi\hookrightarrow E_h^{\otimes k},\qquad
 \eta[P]=\upsilon_h^{\otimes k}P\Bigl(\sum_i s_i\Bigr)1.
 \label{eq:eta}
\end{equation}
For fixed-order monic division
$\chi(S)P(S)=\sum_i h(s_i)Q_i(\mathbf s)$, the ordered primitive
$\sum_i(-1)^{i-1}Q_i\,ds_1\wedge\cdots\widehat{ds_i}\cdots\wedge ds_k$
has image $\sum_i(h_iQ_i-t_iQ_i+u\partial_iQ_i)$.
It follows that
\begin{equation}
 [\chi(S)P(S)]_{\HH_k}
    =\sum_i[t_iQ_i-u\partial_{s_i}Q_i]_{\HH_k}.
 \label{eq:cyclic-defect}
\end{equation}
The quotient by the cyclic image and this parameter defect remain in
the full family.

\section{A genuine finite-field mixed extension}

Choose the same finitely generated coefficient ring $R_0\subset\C$
containing the coefficients of $h$, the matrices of the complete
$\upsilon_h$ and its inverse, and $1/(d+1)!$.  Fix a specified
finite-field point $R_0\to\kappa=\F_Q$ of characteristic $p>d+1$.
For the rest of this section polynomial coefficients mean their images
under this specified map.  Let
\begin{align}
 S&=\Gm_u\times\A^k_{\mathbf t},\qquad
 \pi:\A^k_{\mathbf s}\times S\longrightarrow S,\nonumber\\
 \mathcal L&=\mathcal L_\psi
       \left(\frac{\sum_i\Phi_h(s_i)-\sum_i t_i s_i}{u}\right),
       \qquad \mathcal F=R^k\pi_!\mathcal L.
 \label{eq:finite-family}
\end{align}
Here $\psi$ is a fixed nontrivial additive character with values in the
chosen $\ell$-adic coefficient field, $\ell\ne p$.

\begin{proposition}[The actual open family and boundary stalk]
The sheaf $\mathcal F$ in \eqref{eq:finite-family} is lisse, of rank
$d^k$, and pointwise pure of weight $k$; all $R^j\pi_!\mathcal L$ with
$j\ne k$ vanish.  Fix any finite-field parameter value
$\mathbf t=\mathbf t_0$, after a residue-field extension when necessary,
and write $\mathcal F_{\mathbf t_0}$ for its pullback to $\Gm$.
For $j:\Gm\hookrightarrow\A^1_u$ and $i:\{0\}\hookrightarrow\A^1_u$
there is an exact sequence
\begin{equation}
 0\longrightarrow j_!\mathcal F_{\mathbf t_0}
  \longrightarrow j_*\mathcal F_{\mathbf t_0}
  \longrightarrow i_*(V^I)\longrightarrow0,                \label{eq:boundary-sequence}
\end{equation}
where $V$ is the local geometric generic representation and $I$ its
inertia group.  The boundary stalk $V^I$ is mixed of weights at most $k$.
\end{proposition}
\begin{proof}
The highest homogeneous phase is
$\sum_i s_i^{d+1}/((d+1)u)$.  Its partial derivatives $s_i^d/u$
cannot vanish simultaneously at a projective point.  Therefore the
explicit coefficient map from $S$ to the universal degree-$(d+1)$
polynomial parameter space lands entirely in Deligne's good open.
Weil II, Lemma 3.7.3, proves lissity of the universal proper-support
cohomology sheaves there.  Proper-support base change pulls these sheaves
to \eqref{eq:finite-family}.  Statement 3.7.2.3, with polynomial degree
$d+1$ and affine dimension $k$, proves the asserted concentration, rank,
and weight.  This check uses no simplicity assumption on finite critical
roots.

The stalk of $j_*\mathcal F_{\mathbf t_0}$ at zero is $V^I$; the stalk
of $j_!\mathcal F_{\mathbf t_0}$ there is zero.  On the open the map is
the identity, proving \eqref{eq:boundary-sequence} on every geometric
stalk.  The extension by zero is mixed of weight $k$.

On an open finite-index inertia subgroup with unipotent action, let $N$
be the actual logarithmic monodromy operator and $M$ its monodromy
filtration.  Weil II 1.8.4 applies to the lisse pointwise pure
$\mathcal F_{\mathbf t_0}$ and proves that
$\operatorname{Gr}_r^M V$ is pure of weight $k+r$.
By 1.8.8, $V^I$ is obtained from $\ker N$ by finite-group invariants;
the filtration it inherits has no graded pieces with $r>0$.
Its nonzero graded pieces therefore have weights at most $k$.
No graded rank or numerical matrix for $N$ is inserted in this argument.
\end{proof}

This supplies an actual use of Deligne's mixed-boundary machinery on
the finite-field branch.  The coefficient specialization of
\eqref{eq:lattice}, and the theta source map
\eqref{eq:source-square}, are separately explicit.  A comparison from
their marked complex fibre to the local representation in
\eqref{eq:boundary-sequence} has not been proved by these constructions.
The complex boundary data relevant to such a comparison are calculated
next, with all original coefficients retained.

\section{Exact division and the full boundary connection}

For $0\le b<d$ perform monic division in the original variable $s$:
\begin{equation}
 f_t(s)s^b=(h(s)-t)q_b(s,t)+r_b(s,t),\qquad
 \deg_s r_b<d,\quad \deg_s q_b\le b+1\le d.
 \label{eq:division}
\end{equation}
Let $C_h(t)$ and $B_h(t)$ be the matrices whose $b$th columns are the
coefficients of $r_b$ and $\partial_s q_b$, respectively, in the
ascending monomial basis $1,s,\ldots,s^{d-1}$.

\begin{lemma}[Exact differential remainder]
In the one-factor top lattice,
\begin{equation}
 [f_t s^b]=[r_b-u\partial_s q_b].
 \label{eq:reduction}
\end{equation}
There is no further reduction of the displayed quotient derivative.
\end{lemma}
\begin{proof}
Subtract $Lq_b=(h-t)q_b+u\partial_s q_b$ from
\eqref{eq:division}.  The derivative has degree below $d$, so the
right side is already in the retained remainder basis.
\end{proof}

\begin{theorem}[Degree-correct flat connection]
On the one-factor complex localized at $u$, the degree-one and
degree-zero connections in the $u$ direction are
\begin{equation}
 \nabla_u^1=\partial_u-\frac{f_t}{u^2},\qquad
 \nabla_u^0=\partial_u-\frac{f_t}{u^2}+\frac1u.
 \label{eq:chain-one}
\end{equation}
The induced top connection in the original monomial lattice is exactly
\begin{equation}
 \boxed{\nabla_u=\partial_u-\frac{C_h(t)}{u^2}
                                  +\frac{B_h(t)}u.}
 \label{eq:connection-one}
\end{equation}
On the ordered $k$-factor complex the degree-$j$ operator and its top
cohomology operator are
\begin{align}
 \nabla_u^{(j)}&=\partial_u-
       \frac{\sum_i f_{t_i}(s_i)}{u^2}+\frac{k-j}{u},
       \label{eq:chain-all}\\
 \boxed{\nabla_u}&=\boxed{\partial_u-
       \frac{\sum_i C_{h,i}(t_i)}{u^2}
                +\frac{\sum_i B_{h,i}(t_i)}u.}
       \label{eq:connection-all}
\end{align}
These commute with the existing parameter connections
$\nabla_{t_i}=\partial_{t_i}-s_i/u$ and are compatible with the
original complex differential.
\end{theorem}
\begin{proof}
The commutators on polynomial coefficients are
\[
 [L,\partial_u]=-\partial_s,\qquad
 [L,-f_t/u^2]=-(h-t)/u,
\]
so $[L,\nabla_u^1]=-L/u$ and
$\nabla_u^1L=L(\nabla_u^1+1/u)=L\nabla_u^0$.
This proves the chain identity, including the degree-zero correction.
The differential remainder \eqref{eq:reduction} proves
\eqref{eq:connection-one}.

For \eqref{eq:chain-all}, the commutator with each $L_i$ is
$-L_i/u$, and the scalar difference between adjacent exterior degrees
contributes $+L_i/u$.  Their cancellation proves chain compatibility
in the original order.  Top cohomology and
\eqref{eq:reduction} give \eqref{eq:connection-all}.
Finally, differentiating $-s_i/u$ in $u$ gives $+s_i/u^2$, whereas
the commutator of $-\sum f_{t_l}/u^2$ with $\partial_{t_i}$ is
$-s_i/u^2$.  All other terms commute.  This proves the asserted
flatness.
\end{proof}

\begin{proposition}[The exact first-order trace]
$B_h(t)$ is independent of $t$.  In the unchanged ascending monomial
basis it is upper triangular and satisfies
\begin{equation}
 (B_h)_{bb}=\frac{b+1}{d+1},\qquad
 \boxed{\Tr B_h=\frac d2.}
 \label{eq:trace}
\end{equation}
\end{proposition}
\begin{proof}
In \eqref{eq:division}, compare coefficients of degrees greater than
$d$.  The perturbation $-ts^{b+1}$ and the product $-tq_b$ both have
degree at most $d$.  All coefficients of $q_b$ except possibly its
constant coefficient are therefore determined independently of $t$.
Differentiation removes that constant.  The leading term of $q_b$ is
$s^{b+1}/(d+1)$, so its derivative has degree $b$ and leading coefficient
$(b+1)/(d+1)$.  Summing those diagonal entries proves
\eqref{eq:trace}.
\end{proof}

The matrix $B_h$ alone is not asserted to be the local monodromy
residue: the order-two coefficient $C_h(t)$ is present and generally
does not commute with $B_h$.  Equations \eqref{eq:connection-one} and
\eqref{eq:connection-all} retain their exact relationship.

\section{The full primary jets at the boundary}

\begin{proposition}[Leading coefficient on every original primary block]
Retain the Chinese-remainder summands
$E_\rho=\C[s]/(s-\rho)^{m_\rho}$ and their projectors.  At $t=0$,
\begin{equation}
 C_h(0)|_{E_\rho}=\Phi_h(\rho)I_{E_\rho},\qquad
 A_h|_{E_\rho}=\rho I_{E_\rho}+N_\rho,\qquad
 \dim E_\rho=m_\rho.
 \label{eq:critical-block}
\end{equation}
For the ordered tuple $(\rho_1,\ldots,\rho_k)$ the leading matrix in
\eqref{eq:connection-all} is scalar
$\sum_i\Phi_h(\rho_i)$ on a block of dimension
$\prod_i m_{\rho_i}$.
\end{proposition}
\begin{proof}
Write $h(s)=(s-\rho)^{m_\rho}a_\rho(s)$, where
$a_\rho(\rho)\ne0$.  Integration of its Taylor polynomial, with the
original constant retained, gives the exact identity
\begin{equation}
 \Phi_h(s)-\Phi_h(\rho)=
 \sum_{j\ge0}\frac{a_\rho^{(j)}(\rho)}{j!}
       \frac{(s-\rho)^{m_\rho+j+1}}{m_\rho+j+1}.
 \label{eq:full-taylor}
\end{equation}
The sum is finite.  Every term is divisible by
$(s-\rho)^{m_\rho+1}$.  The matrix $C_h(0)$ is multiplication by
$[\Phi_h]_h$, by its definition in \eqref{eq:division}; hence its
restriction is the first expression in \eqref{eq:critical-block}.
Multiplication by $s$ remains $\rho+N_\rho$ on that same full module.
The tensor statement follows by applying the sum of these matrices
to the ordered tensor projectors.
\end{proof}

Equal sums of critical values retain their separate original projectors
and source labels.  The scalar exponential factors for two tuples have
the exact modulus ratio
\begin{equation}
 \exp\!\left(\operatorname{Re}
      \frac{\sum_i\Phi_h(\rho_i)-\sum_i\Phi_h(\rho_i')}{u}\right).
 \label{eq:directions}
\end{equation}
For a nonzero difference, equal modulus occurs on
$\arg u=\arg(\text{difference})+\pi/2\pmod\pi$.
These are candidate Stokes directions of the displayed exponential
factors.  Equation \eqref{eq:directions} supplies no assertion that a
particular Stokes entry is nonzero, and does not replace the calculation
of the matrices for \eqref{eq:connection-all}.

\section{Exact ordered period determinant}

Keep the original rays and contours
\begin{equation}
 \vartheta_j=\frac{\arg u+\pi}{d+1}+\frac{2\pi j}{d+1},
 \qquad \Gamma_j=-\ell_0+\ell_j\quad(1\le j\le d),
 \qquad
 (\Pi_h)_{jb}=\int_{\Gamma_j}e^{f_t(s)/u}s^b\,ds.
 \label{eq:period}
\end{equation}
The leading radial exponent is
$-r^{d+1}/((d+1)|u|)$; it proves absolute convergence and
coefficient differentiation on compact parameter sets.  Contours on
nearby parameters deform within the same decaying sectors, with their
original orientations.  Every formula below uses a fixed retained
branch for $u$ and the ordered rays.

\begin{theorem}[Full period determinant, including its constant]
Let $c_h(t)=\Tr C_h(t)$, and put
\begin{align}
 \zeta&=e^{2\pi i/(d+1)},\qquad \beta_b=\frac{b+1}{d+1},
       \nonumber\\
 W_{jb}&=\zeta^{j(b+1)}-1
       \quad(1\le j\le d,\ 0\le b<d),\nonumber\\
 K_d&=(d+1)^{-d/2}e^{\pi i d/2}
           \left(\prod_{b=0}^{d-1}\Gamma(\beta_b)\right)\det W.
       \label{eq:constant}
\end{align}
Then $K_d\ne0$ and
\begin{equation}
 \boxed{\det\Pi_h(u,t)=K_d u^{d/2}e^{c_h(t)/u}.}
 \label{eq:det-one}
\end{equation}
For the ordered external product $\Pi_k=\bigotimes_i\Pi_h(u,t_i)$,
\begin{equation}
 \boxed{\det\Pi_k
   =K_d^{k d^{k-1}}u^{k d^k/2}
         \exp\!\left(\frac{d^{k-1}\sum_i c_h(t_i)}u\right).}
 \label{eq:det-all}
\end{equation}
Both identities retain all lower coefficients of $h$ and extend across
critical-root collisions.
\end{theorem}
\begin{proof}
Differentiating \eqref{eq:period} in $u$ inserts $-f_t/u^2$.
The exact reduction \eqref{eq:reduction} gives
\[
 \partial_u\Pi_h=\Pi_h\left(-\frac{C_h(t)}{u^2}
                                    +\frac{B_h}u\right).
\]
The determinant equation, valid as the exterior-power equation even
before invertibility is known, is
\begin{equation}
 \partial_u\det\Pi_h=
     \left(-\frac{c_h(t)}{u^2}+\frac d{2u}\right)\det\Pi_h.
 \label{eq:det-ode}
\end{equation}
Thus the remaining factor after division by $u^{d/2}e^{c_h(t)/u}$
is independent of $u$ on a contour branch.

We prove its independence of $t$ and all lower coefficients without
deleting their effects.  On the dense open of simple roots of $h-t$,
\[
 c_h(t)=\sum_{h(\rho)=t}f_t(\rho),\qquad
 \partial_t c_h(t)=-\sum_{h(\rho)=t}\rho=-\Tr A_h(t).
\]
The terms containing derivatives of the roots vanish because
$f_t'(\rho)=0$.  Both sides are polynomial expressions in the
coefficients and $t$, so this identity holds across collisions too.
The period equation
$\partial_t\Pi_h=\Pi_h(-A_h(t)/u)$ cancels the derivative of
$e^{c_h(t)/u}$, proving that the remaining factor is independent of $t$.

Differentiation in $h_b$ inserts $s^{b+1}/((b+1)u)$.  For a column
$s^j$, its numerator has degree $j+b+1\le2d-1$.
In division by $h-t$ the quotient derivative has degree at most
$j+b-d<j$.  Consequently the differential-reduction correction has
zero trace.  The determinant logarithmic derivative in $h_b$ is
$\Tr M_{s^{b+1}}/((b+1)u)$.  On the simple-root open this equals
$(\partial_{h_b}c_h(t))/u$, since again the root derivative is multiplied
by $f_t'(\rho)=0$.  Polynomial continuation proves the identity for all
coefficients.  Therefore the remaining factor has zero derivative in
every lower coefficient.

It remains to compute that factor at $h=s^d,t=0$, with the original
contours.  Substitution $v=r^{d+1}/((d+1)|u|)$ on each ray gives
\begin{equation}
 (\Pi_h)_{jb}
  =[(d+1)u]^{\beta_b}e^{\pi i\beta_b}
        \frac{\Gamma(\beta_b)}{d+1}
             (\zeta^{j(b+1)}-1).
 \label{eq:gamma-columns}
\end{equation}
The sum of the $\beta_b$ is $d/2$, so taking determinants gives
exactly \eqref{eq:constant} and \eqref{eq:det-one}.
For invertibility of $W$, suppose
$\sum_{r=1}^d a_r(z^r-1)$ vanishes at the $d$ nontrivial roots of
$z^{d+1}=1$.  It also vanishes at $z=1$.  Its degree is at most $d$,
so it is identically zero and every $a_r=0$.  All gamma factors are
nonzero, proving $K_d\ne0$.

The periods and their parameter derivatives are entire in the lower
coefficients on compact sets by the fixed leading radial decay.
This justifies the coefficient continuation used above and proves the
identities also at collisions.  Finally,
$\det(A\otimes B)=(\det A)^{\dim B}(\det B)^{\dim A}$ in the
retained ordered product bases gives \eqref{eq:det-all}.
\end{proof}

\section{Exact comparison with the original source metric}

The period Gram determinant following from \eqref{eq:det-all} is
\begin{equation}
 \det(\Pi_k^*\Pi_k)=|K_d|^{2k d^{k-1}}|u|^{k d^k}
    \exp\!\left(2d^{k-1}\operatorname{Re}
                    \frac{\sum_i c_h(t_i)}u\right).
 \label{eq:period-gram}
\end{equation}
Retain, independently, the actual one-factor source Gram
\begin{equation}
 (G_\theta)_{ab}=\frac1{2\pi}\int_{\mathbb R}
     \overline{(1/2+iy)}^{a}(1/2+iy)^b
       |v_h(1/2+iy)|^2\,dy,\qquad 0\le a,b<d.
 \label{eq:theta-gram}
\end{equation}
This is the Gram of $\mathcal T_h(s^b)=D^bF_h$ in the original
$L^2(dx)$ source norm by Mellin--Plancherel.  Its complete multiplier is
$v_h=g/h$.  It is positive definite: a zero norm forces the continuous
source function $\mathcal T_h(P)$ to vanish; its Mellin transform
$Pv_h$ is then zero, and the nonzero entire $v_h$ forces $P=0$.
The original ordered tensor source Gram is $G_\theta^{\otimes k}$.

The exact comparison of the two finite forms is the endomorphism
\begin{equation}
 T_{\rm per}=(G_\theta^{\otimes k})^{-1}(\Pi_k^*\Pi_k).
 \label{eq:metric-map}
\end{equation}
It is self-adjoint and positive in the original source metric, since
\[
 \langle v,T_{\rm per}w\rangle_{G_\theta^{\otimes k}}
     =v^*\Pi_k^*\Pi_k w,
 \qquad
 \langle v,T_{\rm per}v\rangle_{G_\theta^{\otimes k}}
     =\|\Pi_kv\|^2>0\quad(v\ne0).
\]
Its determinant is exactly \eqref{eq:period-gram} divided by
$(\det G_\theta)^{k d^{k-1}}$.  This preserves both metrics rather
than identifying them.

On the weighted cyclic image \eqref{eq:eta}, the period form is
$\eta^*\Pi_k^*\Pi_k\eta$.  Its relation to the original cyclic
source metric retains the rectangular map $\Pi_k\eta$, the complement
of $\eta$, and the transverse map
\[
 E\xrightarrow{\eta}E_h^{\otimes k}
 \xrightarrow{\sum_i\mathcal R_{h,i}}E_h^{\otimes k}
 \longrightarrow E_h^{\otimes k}/\eta E.
\]
The full determinant \eqref{eq:period-gram} does not determine the
singular values of this rectangular restriction.  The displayed maps
are the exact relationship; no original relation norm has been set to
zero.

\section{Compatibility with independently retained mixed faces}

For an original coefficient mask $A$, retain its $A$-indexed coefficient
complex and the original outer theta-leg label.  Let
$j_A^B$ be the specified constant zero-insertion linear map for
$A\subseteq B$.  It inserts represented zero in each added slot.
The preceding complex differential and connection act with the same
coefficients in each slot.  Therefore
\begin{equation}
 j_A^B\mathcal D_A=\mathcal D_Bj_A^B,
 \qquad j_A^B\nabla_A=\nabla_Bj_A^B.
 \label{eq:faces}
\end{equation}
These equations hold componentwise on retained coordinates; on the
inserted coordinates each derivative and matrix coefficient sends zero
to zero.  Since $j_A^B$ is constant in all parameters, differentiation
commutes with it as well.

On the carrier retaining its mask, a vector killed by an observation
therefore remains the receiving represented zero with that mask.
The empty coefficient face maps to the zero vector in the receiving
face under zero insertion; this face-group map is not a
$\tau$-preserving pointed lift.  A nonempty outer theta-leg label
remains present when the coefficient mask is empty.  No colimit over
the entire mask poset is taken, since its terminal full face would
discard the independent face labels.

The scalar structural specialization is still
\[
 \mathbb F_{1,\tau}\longrightarrow G(\C[u,\mathbf t])
       \xrightarrow{G(\operatorname{ev}_0)}G(\C).
\]
It sends represented $u,t_i$ to represented zero $e$, and $\tau$ to
$\tau$.  The original absolute point, its source, and the coefficient
specialization are retained.  Equation \eqref{eq:faces} is an explicit
compatibility with the mixed face maps; it supplies no claim that a
Boolean mask itself is a Deligne weight.

\section{What follows in the programme}

The $u$-connection, critical-block coefficients, determinant and metric
map above are completed calculations on the original family.  The
finite-field lisse family and its invariant boundary weight bound are
actual applications of Deligne's machinery.  The next calculations
are the local extensions and Stokes or inertia matrices for this phase,
their maps through the original mixed faces, and the induced original
relation forms on the weighted cyclic image.  The parameter defect
\eqref{eq:cyclic-defect}, the transverse map, and all source
cross-pairings remain in that task.  No uniform four-volume upper bound,
arithmetic Frobenius comparison, or failure theorem for the programme
is asserted here.

\section*{Source and verification record}
The original packet-family inputs are those of
\emph{Tau Marked Product Four Endpoint}, 13 September 2026,
\texttt{NOTE.tex}, read in full, with SHA256
\begin{center}\ttfamily\small
40749e2a9d599a322e7b125eb43ad0878\\
a066a51b61da52cc3bb2ffb94cbd841.
\end{center}
The relevant primary source is Pierre Deligne,
\emph{La conjecture de Weil II}, Publications math\'ematiques de
l'IH\'ES 52 (1980), 137--252: 1.8.4, 1.8.8, 3.3.1--3.3.6,
3.7.2.3, 3.7.3, 6.1.13 and 6.2.1--6.2.6.  The local typed text of
those statements and surrounding proofs was read for this audit.

The companion exact symbolic checker verifies differential reduction,
flatness, $\partial_tB_h=0$, the trace and diagonal of $B_h$, the
critical-value trace derivative, and complete jet divisibility for four
polynomial fixtures of degrees 1 through 4, including repeated roots.
All passed.  These fixtures are explicitly not asserted to be
arithmetic zero packets; the proofs here apply to the original $h$.
\end{document}
